\section{Problema}
\label{subsec:problema}

Un problema (computacional, evidentemente) es un enunciado que especifica qué se quiere obtener, sin decir cómo. Más formalmente, describe qué entradas son válidas para el problema y la salida esperada para cada una de ellas.
\begin{definicion}{Problema (computacional)}{problema}
Descripción de la salida esperada para cada instancia en un conjunto de entradas.
\end{definicion}

E.g., el \concept{problema de ordenamiento} (\textit{sorting problem}, muy famoso):
\begin{itemize}
  \item \textbf{Entrada}: secuencia de \(n\) números \(a_1, a_2,\ldots,a_n\).
  \item \textbf{Salida}: permutación \(a_1', a_2',\ldots,a_n'\) tal que \(a_1' \leq a_2' \leq \ldots \leq a_n'\)
\end{itemize}
Recuérdese que cada permutación de una secuencia es un posible orden de sus elementos. En este caso, la salida esperada es cualquier permutación en la que los números estén en orden no decreciente.

Una \concept{instancia} de un problema es una entrada específica, e.g., \(1, 4, 2, 3, 2\) es una instancia del problema de ordenamiento, en cuyo caso la salida esperada es \(1, 2, 2, 3, 4\).

En ocasiones los problemas están redactados de forma que la entrada y la salida esperada están ocultas en el enunciado. En estas notas ese nunca será el caso, pues no se pretende ejercitar la habilidad de interpretar un problema, solo la de resolverlo.

\setinputbase{fundamentos/problema}
\inputlist{tipos, refs}